% ── Part II, Chapter 2 ────────────────────────────────────────
% The broken arm — constraint shock, forced reorganization
% Saint-Domingue, 1740s

\chapter{The Cut}

\strandmarker{Mackandal}{Saint-Domingue, 1740s}

\lettrine{T}{he} mill does not stop.

This is the first thing he understands afterward, lying on the
ground with his arm caught in what the arm is no longer. The mill
continues its rotation. The sound it makes does not change. The
process that requires the mill to turn does not pause to
accommodate what has just happened to one of its components.

The system has registered the event and continued.

He is aware of this with a clarity that has nothing to do with
thought---it arrives before thought, before pain in the organized
sense, before anything that could be called a response. It is
simply present: the mill turns, the process continues, the system
does not require him to be intact in order to proceed.

What it requires is that the component be removed and replaced.

This is already happening. Hands he recognizes without being able
to name arrive and do what hands do in this situation. Instructions
are given in the language of urgency, which is the same in every
language. He is moved. The mill recedes. The sound of it becomes
something he is no longer inside of.

\section{}

The recovery is long and he is conscious for most of it.

This turns out to be a different kind of education.

Before, he moved through the system. His knowledge of it was
kinetic, accumulated through action, stored in the body's
responses. The map he had built was a map for a body that could
move through the terrain it described. It assumed a particular
set of capabilities and organized itself around them.

Now the body cannot do most of what the map assumed.

He expected this to diminish the map. Instead it clarifies it.

Lying still, he can observe what moving prevented him from seeing.
The system continues around him, and because he is no longer a
moving part of it, he can perceive its motion from something
closer to outside. Not fully outside---he is still within it,
still subject to its logic, still a node even if a temporarily
non-functioning one---but differently positioned than before.

The angles are new.

He watches the way instructions travel. He watches where they
originate and where they arrive and what happens in the space
between. He watches which parts of the system are monitored and
which are assumed to be functioning without verification. He
watches what the system cannot see about itself.

He has more time than he has ever had to watch anything.

\section{}

The arm, when it has finished becoming what it will be, is not
what it was.

He examines this without sentimentality, which is not the same as
without feeling. The feeling is present and he allows it to be
present and he does not permit it to be the only thing that is
present. Alongside it, there is the question of what this changes
and what it does not.

What it changes: the set of actions the body can perform. The
map that assumed certain capabilities must be revised. Some paths
through the system that were previously available are no longer
available. Some tasks that were previously assignable to him will
be assigned elsewhere.

What it does not change: the map itself, the knowledge accumulated
within it, the understanding of where the system is thin and where
it is dense and where its coherence depends on assumptions that
have not been tested. None of this was stored in the arm.

There is a reconfiguration available to him.

Not a consolation---he is not interested in consolations, which
are a way of not thinking about something---but a genuine
structural observation. The system has changed what he can do
within it. It has not changed what he knows about it. And what
he knows about it is, he is beginning to understand, more
dangerous to the system than what he could do within it.

A body that functions normally inside a system is useful to the
system.

A mind that understands the system's limits is something else.

\section{}

They assign him to the livestock.

This is, from the system's perspective, a rational reallocation:
a damaged component redirected to tasks that do not require full
function. The system has processed his injury and produced an
appropriate adjustment. It has, in its terms, solved the problem
he presents.

He accepts the reassignment without comment.

What the system cannot know is that the livestock move. They move
between the plantation and the pastures and the water and back
again, and the person responsible for them moves with them, and
the paths they take pass through or near the edge where the
managed land ends and the unmanaged land begins.

He will have reason now to stand at that edge regularly.

He will have reason to know who else stands there, and when, and
what passes between the managed and the unmanaged in both
directions---what goes in, what comes out, what the forest holds
that the plantation does not know it holds.

The system has solved its problem.

In doing so it has created the conditions for a different kind
of solution entirely---one that the system's logic cannot
anticipate because it cannot see far enough outside itself to
know that anything is being solved.

He begins his new work.

The mill turns behind him, audible at this distance, regular,
indifferent.

He does not look back at it.

